\documentclass{hasel_thesis}
\usepackage[utf8]{inputenc}
\usepackage[T1]{fontenc}
\usepackage{geometry}
\usepackage{graphicx}
\usepackage{float}
\usepackage{hyperref}
\usepackage[english,indonesian]{babel}

\title{Optimasi Pemrosesan Template LaTeX Menggunakan Pendekatan Berbasis AI}
\author{Academic Support Agent}
\date{\today}

\begin{document}
\maketitle

\begin{abstract}

\end{abstract}

\chapter{PENDAHULUAN}


\chapter{TINJAUAN PUSTAKA}


% Original content (commented out):
% \documentclass{hasel_thesis}

\thesisType{Bachelor/Master Thesis}
\date{\today}
\title{Example title}
\subtitle{Example subtitle}
\author{<first name> <last name>}
\home{<geburtsort>} % Birth town
\country{<heimatland>} % Birth country
\legi{xx-xxx-xxx}
\prof{Prof. Dr. Thomas Fritz}
\assistent{Assistant}
\email{email-address}
\repourl{url if available}
\begindate{start date of thesis}
\enddate{end date of thesis}

\begin{document}
\maketitle

\frontmatter

\begin{acknowledgements}
\end{acknowledgements}

\begin{abstract}
\end{abstract}

\begin{zusammenfassung}
\end{zusammenfassung}

\tableofcontents
\listoffigures
\listoftables
\lstlistoflistings

\mainmatter
\chapter{Chapter}
\section{Section}
%
\subsubsection{Subsubsection}
\fig[.5\textwidth]{logo_hasel}{Our logo}{logo}

\subsection{Subsection}
%
\paragraph{Paragraph.} Always with a point, see \citeauthor{texbook}~\cite{texbook} and \citeauthor{latex:companion}~\cite{latex:companion}.

\begin{lstlisting}[caption=An example code snippet]
/**
 * Javadoc comment
 */
public class Foo {
	// line comment
	public void bar(int number) {
		if (number < 0) {
			return; /* block comment */
		}
	}
}
\end{lstlisting}

\printbibliography


\appendix
\chapter{First Appendix}
\chapter{Second Appendix}

\backmatter

\end{document}


\end{document}